%=====================================================================
% Advanced Fiber Materials / Springer Nature LaTeX Manuscript
%=====================================================================

\documentclass[pdflatex,sn-mathphys-num,iicol]{sn-jnl}

%=====================================================================
% Packages
%=====================================================================
\usepackage{graphicx}
\usepackage{subcaption}
\usepackage{multirow}
\usepackage{amsmath,amssymb,amsfonts}
\usepackage{amsthm}
\usepackage{mathrsfs}
\usepackage[title]{appendix}
\usepackage{xcolor}
\usepackage{hyperref}
\usepackage{textcomp}
\usepackage{manyfoot}
\usepackage{booktabs}
\usepackage{algorithm}
\usepackage{algorithmicx}
\usepackage{algpseudocode}
\usepackage{url}
\usepackage{lmodern}
\usepackage[T1]{fontenc}
\usepackage{tikz}
\usepackage{hyperref}
\usepackage{xurl}


%=====================================================================
% Theorem settings
%=====================================================================
\theoremstyle{thmstyleone}
\newtheorem{theorem}{Theorem}

\theoremstyle{thmstyletwo}
\newtheorem{example}{Example}
\newtheorem{remark}{Remark}

\theoremstyle{thmstylethree}
\newtheorem{definition}{Definition}

\begin{document}

%=====================================================================
% TITLE
%=====================================================================

\title[SEM Defect Quantification for Fiber-Related Perovskite Photovoltaics]
{Interpretable SEM-Based Microdefect Quantification for Fiber and Flexible Perovskite Photovoltaic Materials}

% AUTHORS
%=====================================================================

\author*[1]{\fnm{Sahil} \sur{Soni}\thanks{ORCID: \href{https://orcid.org/0009-0005-0499-1694}{0009-0005-0499-1694}}}
\email{sahilsonii369@gmail.com}

\author[1]{\fnm{Dr. Archana Y.} \sur{Chaudhari}}
\email{archana.chaudhari@sitpune.edu.in}

\affil[1]{
  \orgname{Symbiosis Institute of Technology, Pune Campus,
  Symbiosis International (Deemed University)},
  \orgaddress{\city{Pune},
  \state{Maharashtra},
  \country{India}}
}


%=====================================================================
% ABSTRACT
%=====================================================================

\abstract{
Fiber-related and flexible perovskite photovoltaic materials are
promising for lightweight, wearable, and green-energy harvesting
applications, but their device reliability depends strongly on
thin-film morphology. Microstructural defects such as pinholes, voids,
incomplete coverage, and non-uniform grains can create leakage
pathways, enhance charge recombination, reduce effective
light-absorbing area, and accelerate instability. Recent experimental
studies have shown that nanovoids of approximately 100--200~nm
diameter, directly observable in Scanning Electron Microscopy (SEM)
images, are associated with increased trap-state density and
measurably reduced fill factor and power conversion efficiency in
wide-bandgap perovskite devices. Eliminating such voids through
process optimisation has been shown to improve certified efficiency
from approximately 22\% to over 23\% in state-of-the-art inverted
devices. Manual SEM inspection, however, is slow, subjective, and
difficult to scale for fabrication quality control. This study presents
an interpretable computer-vision pipeline, released as the open-source
Python library \textbf{MicroDefectCV}, for automated SEM-based
microdefect quantification in perovskite films. Operating without
deep learning or labelled data, the package implements morphology-aware
detection modes to quantify PbI$_2$ excess, dark pinholes, and
grain-boundary defects. The method outputs binary defect masks,
YOLO-format bounding-box annotations, defect count, and defect-area
percentage, enabling direct connection between SEM morphology and
material-quality indicators relevant to photovoltaic device performance.
The classical pipeline is benchmarked against manual annotation and
six deep-learning YOLO models on a 4{,}701-patch dataset. While the
optimal YOLO11m model requires significant GPU training and labelling
effort, MicroDefectCV provides a lightweight, reproducible, and
annotation-free route for scalable SEM morphology analysis in
perovskite material development.
}

%=====================================================================
% KEYWORDS
%=====================================================================

\keywords{
Perovskite photovoltaics, SEM defect detection, MicroDefectCV,
OpenCV adaptive pipeline, microdefect quantification, Python package,
fiber solar cells, flexible perovskite, YOLO benchmarking
}

\maketitle

%=====================================================================
% 1. INTRODUCTION
%=====================================================================

\section{Introduction}\label{sec:introduction}

\begin{figure*}[!htbp]
\centering
\IfFileExists{Fig0_perovskite_concept.png}{
  \includegraphics[width=0.80\textwidth]{Fig0_perovskite_concept.png}
}{
  \fbox{\begin{minipage}[c][5cm][c]{0.78\textwidth}\centering
  Place \texttt{Fig0\_perovskite\_concept.jpg} in the directory.\end{minipage}}
}
\caption{Layered architecture of a flexible perovskite solar cell (Au / HTL / Perovskite / ETL / ITO / Flexible Substrate).}
\label{fig:perovskite_concept}
\end{figure*}

Perovskite photovoltaic materials have emerged as a leading
thin-film semiconductor class for next-generation solar-energy
conversion, combining high optical absorption coefficients
($\sim$10$^5$~cm$^{-1}$), broadly tunable band gaps
(1.2--3.0~eV via halide and cation substitution),
low defect-state energies, and solution-processable fabrication routes
compatible with both rigid and flexible substrates
\cite{kojima2009,snaith2013,green2014,park2015}.
Recent progress is striking: certified efficiencies now exceed 27\%
for narrow-bandgap inverted devices \cite{green2025}, while
wide-bandgap (1.65--1.80~eV) formamidinium-based compositions have
reached 23.42\% with fill factors of 86.8\% \cite{chang2026}.
Perovskite/silicon tandem and triple-junction architectures have
achieved 30.97\% and 27.06\% respectively \cite{chang2026,zheng2025}.

Despite these achievements, thin-film morphology remains the
primary determinant of device reliability and batch-to-batch
reproducibility. During perovskite film formation, the high-boiling-point
solvent DMSO becomes trapped within the crystallising film and escapes
during thermal annealing, leaving behind nanovoids of
$144.6 \pm 56.9$~nm average diameter at the perovskite/HTL buried
interface \cite{chang2026,chen2021}. These voids increase
non-radiative recombination, raise trap-state density by 14.8\%,
and directly degrade fill factor and power conversion efficiency
\cite{chang2026}. Pinholes, incomplete film coverage, and non-uniform
grain crystallisation similarly create leakage channels and
photo-degradation pathways. For fiber-shaped and flexible
perovskite solar cells as depict in figure \ref{fig:perovskite_concept}, where deposition occurs on curved or
mechanically deformable surfaces, these morphological failure modes
are further amplified by substrate curvature, solvent-evaporation
anisotropy, and mechanical strain
\cite{balilonda2019,hu2016,qiu2016,diGiacomo2016}.

Scanning Electron Microscopy (SEM) is the standard characterisation
technique for thin-film morphology assessment at the nanoscale.
Top-view and cross-sectional SEM images reveal pinhole distribution,
grain boundary character, void density, incomplete film coverage,
and interfacial delamination \cite{chang2026,zheng2025}. However,
manual SEM inspection is slow, observer-dependent, and cannot scale
to the image volumes required for fabrication quality control in a
high-throughput manufacturing environment.

This work presents \textbf{MicroDefectCV}~\cite{microdefectcv2025},
an open-source Python library (\texttt{pip install microdefectcv})
that converts raw perovskite SEM images into quantitative defect
descriptors without any deep learning, labelled training data, or
GPU resources. The pipeline implements six morphology-aware detection
modes and an automatic image-type classifier based on intensity
statistics, combining CLAHE contrast enhancement, percentile-based
dynamic thresholding, grain-boundary suppression, top-hat transforms
for bright-particle detection, and contour-based shape filtering.
Output is provided as binary defect masks, YOLO-format bounding-box
labels, defect count, and defect-area percentage --- material-quality
indicators directly comparable across fabrication batches and
processing conditions. To characterise the annotation landscape and
benchmark MicroDefectCV against supervised approaches, three
annotation strategies (zero-shot AutoDistill, standalone SAM, and
manual annotation) were evaluated, followed by a six-model YOLO
benchmark on the resulting 4{,}701-patch dataset.

fiber-integrated devices, where film morphology may vary along curved
or elongated geometries. For scalable green-energy manufacturing,
there is a need for automated and interpretable tools that can convert
SEM morphology into quantitative defect indicators.

Traditional computer-vision methods such as global thresholding,
Otsu thresholding, edge detection, and watershed segmentation are
commonly used for image segmentation
\cite{otsu1979,canny1986,vincent1991}. These techniques are simple
and computationally efficient, but their performance may degrade when
SEM images contain uneven illumination, low contrast, charging
artefacts, or complex nanoscale texture. Deep learning approaches
can improve detection accuracy but require carefully annotated
datasets and high computational resources. For early-stage materials
research with limited annotations, interpretable computer-vision
pipelines remain valuable.

This study proposes an interpretable SEM-based microdefect
quantification pipeline for perovskite photovoltaic materials, with
emphasis on fiber-related and flexible device inspection. The pipeline
combines contrast-limited adaptive histogram equalization (CLAHE),
dynamic thresholding, morphological refinement, and contour-based
filtering to identify pinhole-like and void-like regions of the type
documented in recent high-performance perovskite device studies
\cite{chang2026,zheng2025}.

The main contributions of this work are:

\begin{itemize}
\item \textbf{MicroDefectCV package}: a publicly released, pip-installable
Python library~\cite{microdefectcv2025} providing six morphology-aware
detection modes, auto image-type classification, grain-boundary
suppression, and direct YOLO-format label export --- all without
deep learning or labelled data.
\item \textbf{Annotation-free defect quantification}: the pipeline
automatically detects pinholes, voids, PbI$_2$ bright particles, and
needle crystals from raw SEM images and produces measurable descriptors
(defect count, area fraction, mean diameter) in under one second per
image on CPU.
\item \textbf{Experimental validation through annotation and YOLO benchmarks}:
to motivate and contextualise the classical approach, zero-shot
auto-annotation (AutoDistill), manual annotation, and a six-model
YOLO benchmark (4{,}701 patches, five classes) are reported. Despite
being annotation-free, MicroDefectCV outperforms the YOLO models in
practical throughput, requiring no GPU or labelling effort.
\item \textbf{Materials-science grounding}: defect descriptors are
interpreted through the physical mechanisms linking SEM-observable
voids to trap-state density and certified efficiency, consistent with
recent high-performance perovskite device literature
\cite{chang2026,zheng2025}.
\end{itemize}

%=====================================================================
% 2. MATERIALS RELEVANCE AND SEM MORPHOLOGY MOTIVATION
%=====================================================================

\section{Materials Relevance and SEM Morphology Motivation}
\label{sec:materials_relevance}

Perovskite thin films are the photoelectric active layer in both
rigid and flexible solar cells. Their performance is intimately
connected to film morphology, and SEM imaging is the standard
technique for visualising the nanoscale features that determine
device quality. This section summarises the physical mechanisms
that motivate automated SEM defect quantification.

\subsection{Microdefects in Perovskite Photovoltaic Films}
\label{subsec:microdefects}

In perovskite photovoltaic films, microdefects are closely related to
film formation, crystallisation, solvent evaporation, substrate
wetting, and post-treatment conditions. A visually small defect in an
SEM image may correspond to incomplete perovskite coverage, a local
void, or a weakly connected grain region. Such defects can disturb
charge transport and reduce the spatial uniformity of photocurrent
generation.

Recent experimental evidence provides quantitative context for this
connection. Chang et al.\ \cite{chang2026} reported that SEM images
of wide-bandgap ($\sim$1.67~eV) formamidinium-based perovskite films
deposited on single-component hole-transport layers revealed nanovoids
with an average diameter of $144.6 \pm 56.9$~nm. These voids were
attributed to the segregation of the high-boiling-point solvent DMSO,
which becomes trapped within the perovskite film during deposition and
escapes from the bottom, leaving voids behind
\cite{chang2026,chen2021}. This mechanism is especially relevant
during multi-step spin-coating processes used in both planar and
flexible device fabrication.

When such voids were eliminated through hole-transport layer
engineering, the device power conversion efficiency improved from
21.92\% to 23.04\%, accompanied by improvements in V$_{oc}$, FF,
and J$_{sc}$ \cite{chang2026}. The fill factor reached a certified
value of 86.8\% in the void-free configuration. This demonstrates
that even nanoscale voids visible in SEM images carry measurable
consequences for macroscopic device performance.

Pinholes and voids are particularly important because they can form
electrical leakage channels between the electron- and hole-transport
layers. In a solar cell, this may increase shunt pathways and reduce
open-circuit voltage or fill factor. In flexible and fiber-shaped
perovskite devices, such defects may become more severe because the
active layer is deposited on curved, rough, or mechanically
deformable surfaces. Hence, automated pinhole and void quantification
can support process optimisation for coating uniformity and device
reproducibility.

\subsection{Physical Mechanisms Linking SEM Morphology to Device Performance}
\label{subsec:physical_mechanisms}

Understanding the physical origin of SEM-visible defects is essential
for interpreting automated quantification results as materials-quality
indicators. The following mechanisms connect SEM-observable
morphology to device physics.

\textbf{Void formation and solvent trapping.}
DMSO, commonly used with DMF to dissolve perovskite precursors, tends
to become trapped within the film during spin-coating and escapes from
the bottom during annealing, leaving voids behind
\cite{chang2026,chen2021}. These voids appear as dark regions in
top-view SEM images, making them detectable by intensity-based
segmentation.

\textbf{Non-radiative recombination and trap-state density.}
Trap-state densities at the perovskite/hole-transport-layer interface
have been measured by space-charge-limited current analysis. A
trap-state density of $2.62 \times 10^{16}$~cm$^{-3}$ was reported
for void-containing films, reduced by 14.8\% to
$2.24 \times 10^{16}$~cm$^{-3}$ when voids were eliminated
\cite{chang2026}. This quantitative link between SEM-observable
morphology and trap density is the physical justification for
treating the pipeline output as a materials-quality indicator.

\textbf{Phase segregation and illumination stability.}
Perovskite films with higher void coverage show more pronounced
photo-induced halide segregation under continuous illumination,
as monitored by time-resolved photoluminescence \cite{chang2026}.
Charge accumulation at defect-rich regions promotes ion migration
and bandgap redshift, which reduces long-term photovoltaic stability.

\textbf{Interface quality in multi-junction architectures.}
For tandem and triple-junction devices, cross-sectional SEM imaging
is used to confirm layer continuity and absence of voids at buried
interfaces \cite{zheng2025}. In a triple-junction
perovskite/perovskite/silicon device achieving 27.06\% certified
efficiency, cross-sectional SEM confirmed continuous layers without
voids at each subcell interface \cite{zheng2025}. This illustrates
that SEM morphology verification is standard practice even in the
most advanced photovoltaic architectures.

\subsection{Relevance to Fiber-Related and Flexible Photovoltaics}
\label{subsec:fiber_relevance}

Fiber-related and flexible perovskite photovoltaic devices require
uniform deposition over non-traditional geometries such as curved
fibers, flexible polymer substrates, yarn-like structures, or
textile-compatible surfaces. Compared with rigid planar substrates,
these geometries introduce additional challenges in coating thickness
control, surface wetting, mechanical strain, and interfacial adhesion
\cite{balilonda2019,hu2016,qiu2016,diGiacomo2016}.

The void and pinhole formation mechanisms identified for planar
devices are expected to be more pronounced in fiber-shaped and
flexible geometries, where substrate curvature affects solvent
evaporation dynamics and film crystallisation uniformity. The
DMSO-trapping mechanism that produces voids in flat-substrate films
\cite{chang2026} would be further modulated by local substrate
geometry and coating speed on a curved fiber surface. Furthermore,
the void diameter range identified in recent literature
($\sim$100--200~nm \cite{chang2026}) falls within the spatial
resolution range accessible to SEM imaging, making automated
SEM-based detection a practical approach for fiber and flexible
device quality assessment.

Although the SEM images analysed in this study are obtained from
planar or locally flat regions, the proposed quantification strategy
is directly relevant to flexible and fiber-related perovskite
photovoltaics because the same defect types---pinholes, voids, and
incomplete coverage---are critical during scalable coating and device
integration. The method can be extended to curved-fiber SEM images,
flexible substrate films, and roll-to-roll processed perovskite
layers.

\subsection{From Image Objects to Material-Quality Indicators}
\label{subsec:image_to_material}

The proposed method treats detected regions as material-quality
indicators rather than isolated image objects. Each detected defect
is converted into measurable descriptors such as area, equivalent
diameter, spatial density, and defect-area percentage. These values
can be compared across fabrication batches, coating recipes,
annealing temperatures, precursor concentrations, or substrate
treatments.

The average nanovoid diameter of $144.6 \pm 56.9$~nm reported by
Chang et al.\ \cite{chang2026} represents the type of quantitative
descriptor that the proposed pipeline aims to produce automatically
from SEM images, rather than through manual measurement. Such
automated extraction enables larger-scale morphology studies across
many samples and conditions, which is important for process
optimisation in green-energy manufacturing.

The central objective of this work is therefore to bridge SEM image
analysis with materials characterisation by converting visual
morphology into reproducible quantitative indicators. These
indicators are grounded in the physical mechanisms connecting
SEM-observable defects to photovoltaic device performance, as
established by recent experimental literature
\cite{chang2026,zheng2025,hu2023}.

%=====================================================================
% 3. RELATED WORK
%=====================================================================

\section{Related Work}\label{sec:related_work}

The development of automated SEM analysis tools for perovskite
morphology spans classical image processing, deep learning
baselines, and materials characterisation. This section reviews the
most relevant prior work and identifies the gap that MicroDefectCV
is designed to fill.

\subsection{Perovskite and Fiber-Related Photovoltaic Materials}
\label{subsec:psc_related_work}

Perovskite solar cells have progressed rapidly since their early
demonstration as visible-light sensitisers \cite{kojima2009}. Their
high absorption coefficient and compatibility with thin-film
processing make them attractive for flexible and lightweight
photovoltaic systems \cite{snaith2013,green2014,park2015}. Inverted
(p-i-n) device architectures now achieve certified efficiencies
exceeding 27\% for narrow-bandgap compositions \cite{green2025},
and wide-bandgap (1.67--1.73~eV) inverted devices have reached
certified efficiencies of 23.42\% with a fill factor of 86.8\%
\cite{chang2026}. These advances highlight the critical role of
interface quality, film morphology, and defect density in determining
photovoltaic performance.

At the system level, perovskite cells are increasingly integrated
into tandem and triple-junction architectures. A four-terminal
perovskite/silicon tandem achieved 30.97\% efficiency
\cite{chang2026}, while a monolithic perovskite/perovskite/silicon
triple junction reached a certified 27.06\% at 1~cm\textsuperscript{2}
and 23.3\% at 16~cm\textsuperscript{2} \cite{zheng2025}. These
state-of-the-art demonstrations routinely use SEM imaging for
morphological verification, confirming that SEM-based quality control
is a standard tool in high-performance photovoltaic research
\cite{chang2026,zheng2025}.

Fiber-shaped and flexible perovskite solar cells extend these
advantages toward wearable and textile-compatible energy harvesting,
but they also introduce additional challenges in morphology control,
coating uniformity, and mechanical durability
\cite{balilonda2019,hu2016,qiu2016,diGiacomo2016}.

\subsection{SEM Morphology and Defect-Performance Correlation}
\label{subsec:sem_morphology_correlation}

A growing body of experimental work directly links SEM-observable
morphological features to perovskite device performance. Chang et al.\
\cite{chang2026} demonstrated that nanovoids of average diameter
$144.6 \pm 56.9$~nm, observed in top-view SEM images of 1.67~eV
perovskite films, contributed to a trap-state density of
$2.62 \times 10^{16}$~cm$^{-3}$ and enhanced non-radiative
recombination at the buried interface. Eliminating these voids
reduced the trap-state density by 14.8\% and improved the certified
fill factor to 86.8\%.

Hu et al.\ \cite{hu2023} demonstrated a void-free buried interface
for scalable FAPbI\textsubscript{3} perovskite solar modules, showing
that void elimination at the buried interface directly improved module
stability and performance. Chen et al.\ \cite{chen2021} showed that
interfacial defects originating from void-related morphology
degradation were a primary cause of performance loss in perovskite
modules under operational conditions. In high-performance
triple-junction devices, cross-sectional SEM imaging confirmed
void-free and delamination-free interfaces as prerequisites for
achieving high fill factors \cite{zheng2025}.

These results collectively establish that SEM morphology analysis is
not a supplementary characterisation step but a primary method for
assessing perovskite film quality. However, all of the above studies
relied on manual SEM inspection. Automated and quantitative pipelines
that can extract defect descriptors from SEM images with
reproducibility and throughput remain an open need in the field.

\subsection{Computer Vision and Learning-Based Defect Detection}
\label{subsec:cv_related_work}

Classical segmentation methods such as Otsu thresholding, Canny edge
detection, and watershed segmentation are widely used because of
their simplicity and interpretability
\cite{otsu1979,canny1986,vincent1991}. CLAHE is useful for local
contrast enhancement and can improve defect visibility in
low-contrast images \cite{zuiderveld1994}. Deep learning methods such
as U-Net and YOLO have achieved strong performance in segmentation
and object detection tasks \cite{ronneberger2015,redmon2016}.
However, deep learning usually requires larger labelled datasets and
careful generalisation testing. For early-stage materials research
with limited annotations, interpretable computer-vision pipelines
remain valuable. Furthermore, the nanovoid diameter range observed
in perovskite SEM images ($\sim$100--200~nm \cite{chang2026}) is at
a spatial scale where intensity-based segmentation methods are
well-suited, provided sufficient contrast enhancement is applied
before thresholding. Table \ref{tab:literature_review} describe the   literature and research gap of existing studies.

%=====================================================================
% Literature Review Table
%=====================================================================

\begin{table*}[!htbp]
\caption{Representative literature and research gap motivating the present study}
\label{tab:literature_review}
\begin{tabular}{@{}p{2.8cm}p{3.0cm}p{4.6cm}p{4.8cm}@{}}
\toprule
Reference & Area & Key contribution & Gap addressed in this work \\
\midrule
Kojima et al.\ \cite{kojima2009}
  & Perovskite photovoltaics
  & Demonstrated organometal halide perovskites for photovoltaic applications
  & Does not address automated morphology quantification \\[4pt]
Green et al.\ \cite{green2014}
  & Perovskite solar cells
  & Discussed emergence and rapid progress of perovskite solar cells
  & Focus is device progress rather than SEM defect analysis \\[4pt]
Chang et al.\ \cite{chang2026}
  & Wide-bandgap inverted PSC; tandem
  & SEM-visible nanovoids ($144.6 \pm 56.9$~nm) linked to 14.8\% higher
    trap density; void elimination improved certified PCE to 23.42\%
    with FF 86.8\%
  & Manual SEM used; no automated quantification pipeline \\[4pt]
Zheng et al.\ \cite{zheng2025}
  & Perovskite
  /perovskite
  /Si triple junction
  & SEM cross-sections verified void-free interfaces in 27.06\% cells;
    SEM used as standard quality gate
  & SEM used qualitatively; no defect descriptors extracted \\[4pt]
Balilonda et al.\ \cite{balilonda2019}
  & Perovskite solar fibers
  & Reviewed status and challenges of perovskite solar fibers
  & Highlights fiber relevance but not automated SEM quantification \\[4pt]
Hu et al.\ \cite{hu2016}
  & Fiber-shaped perovskite solar cells
  & Reported fiber-shaped perovskite solar-cell architecture
  & Does not provide automated microdefect detection workflow \\[4pt]
Otsu \cite{otsu1979}
  & Thresholding
  & Automatic threshold selection from gray-level histograms
  & Sensitive to SEM contrast and illumination variation \\[4pt]
Canny \cite{canny1986}
  & Edge detection
  & Widely used edge detector
  & Produces edges rather than material-defect regions \\[4pt]
Vincent and Soille \cite{vincent1991}
  & Watershed segmentation
  & Watershed segmentation for digital spaces
  & Can over-segment noisy SEM morphology \\[4pt]
Ronneberger et al.\ \cite{ronneberger2015}
  & Deep segmentation
  & U-Net for image segmentation
  & Requires annotated training data \\[4pt]
Present study
  & SEM defect quantification
  & Interpretable CV pipeline grounded in experimental perovskite literature
  & Defect masks, area fraction, mean defect diameter, baseline validation \\
\botrule
\end{tabular}
\end{table*}

%=====================================================================
% 4. MATERIALS AND DATASET
%=====================================================================

\section{Materials and Dataset}\label{sec:materials_dataset}

The experimental workflow consisted of perovskite film preparation,
structural and optical characterisation, SEM morphology imaging, and
automated defect quantification. The SEM images obtained from the
fabricated perovskite films were used as the primary input for the
proposed computer-vision pipeline.

\subsection{Chemicals and Materials}\label{subsec:chemicals_materials}

Dimethylformamide (DMF, 99.8\%), dimethyl sulfoxide (DMSO, 99.8\%),
ethyl acetate (anhydrous, 99.5\%), hexane (anhydrous, 99.5\%),
acetonitrile (99.8\%), ethanol (anhydrous, 99\%), and 2-propanol
(IPA, anhydrous, 99.9\%) were sourced from Sigma Aldrich. Lead(II)
iodide (PbI$_2$, 99.99\%) was obtained from Tokyo Chemical Industry.
Formamidinium iodide (FAI, $>$98\%), methylammonium bromide (MABr,
$>$98\%), phenethylammonium iodide (PEAI, $>$99\%), n-dodecylammonium
iodide (DDAI, $>$99\%), cyclohexylmethylammonium iodide (CMAI,
$>$99\%), 2-thiophenemethylammonium iodide (TMAI, $>$99\%), and
methylammonium chloride (MACl, $>$98\%) were procured from GreatCell
Solar and used as received \cite{ref24}.

\subsection{Fabrication of Perovskite Films}\label{subsec:fabrication}

ITO-coated or plain glass substrates were cleaned by sequential
sonication in acetone, ethanol, and 2-propanol, followed by 15~min
UV--ozone treatment. High-purity $\delta$-FAPbI$_3$ powder was
synthesised using a one-pot method \cite{ref24}. The powder was
washed three times with ethyl acetate and acetonitrile, dried in a
vacuum oven for more than 24~h, and stored in an N$_2$-filled
glovebox.

For pristine $\alpha$-FAPbI$_3$ films, a 1.8~M precursor solution
was prepared by dissolving 1139~mg $\delta$-FAPbI$_3$ and 40~mg
MACl in 1~mL DMF:DMSO (9:1 v/v). The solution was spin-coated on
the substrates in two steps: 1000~rpm for 10~s and 5000~rpm for 30~s.
During the second step, 100~$\mu$L of ethyl acetate:hexane (7:3 v/v)
antisolvent was applied at 10~s intervals under 30--40\% relative
humidity. The films were annealed at 150$^\circ$C for 20~min to form
$\alpha$-FAPbI$_3$. It should be noted that DMSO, used as part of
the precursor solvent system, has been identified as a key contributor
to void formation in perovskite films: DMSO tends to become trapped
within the film during deposition and escapes during annealing,
leaving voids behind \cite{chang2026,chen2021}. The presence and
size distribution of such voids in the final film were assessed using
SEM imaging and the automated pipeline described in this work.

For two-dimensional (2D) perovskite layers on three-dimensional (3D)
perovskite films, 4~mM solutions of PEAI, TMAI, CMAI, or DDAI in IPA
were sprayed at 3.0~mL/min using N$_2$ gas .

\subsection{Film Characterisation}\label{subsec:film_characterization}

Crystallographic properties were analysed using a Rigaku MiniFlex
X-ray diffractometer (XRD) with Cu K$\alpha$ radiation
($\lambda = 1.5405$~\AA). Absorption spectra were measured using a
Shimadzu UV-3600 Plus UV--vis--NIR spectrophotometer. Steady-state
photoluminescence (PL) was recorded using an Edinburgh Instruments
FS5 spectrofluorometer. Surface morphology was examined using a
JEOL/JSM-F100 scanning electron microscope.

The SEM images were used for morphology-based microdefect
quantification. Pinhole-like and void-like regions observed in SEM
images were treated as morphology indicators associated with
incomplete film coverage, local defect formation, and possible
leakage pathways in perovskite photovoltaic devices
\cite{chang2026,hu2023}. This interpretation is grounded in
experimental literature demonstrating direct correlation between
SEM-observable void density and trap-state density
\cite{chang2026}.

% TODO: Add exact SEM imaging details:
% - accelerating voltage
% - working distance
% - magnification range
% - image resolution (pixels)
% - scale bar value
% - number of SEM images in dataset
% - number of samples

\subsection{SEM Image Dataset for Defect Quantification}
\label{subsec:sem_dataset}

The SEM image dataset was constructed from morphology images obtained
during film characterisation and grouped into five defect and
morphology categories. In total, 440 SEM images were collected,
covering PbI$_2$ excess bright-particle regions, 3D perovskite
pinholes, 3D-2D mixed perovskite pinholes, 3D perovskite background
regions, and 3D-2D mixed perovskite background regions. For the YOLO
training and evaluation pipeline, these 440 base images were split
into 224$\times$224~pixel patches using a sliding-window strategy with
25\% overlap, producing 4{,}701 annotated patches distributed across
training (80\%), validation (10\%), and test (10\%) sets.
Table ~\ref{tab:dataset_summary} summarises the dataset structure.

% TODO: Add exact SEM magnification, accelerating voltage, working
% distance, and image pixel resolution when available.

Manual annotations were prepared to provide a reference for validating
the automated segmentation results. Defect regions were marked based
on visual morphology, local darkness, boundary structure, and material
relevance. Very small isolated dark pixels were not considered defects
unless they satisfied the minimum area and shape criteria. The
expected void size range of $\sim$100--200~nm, as reported for similar
perovskite films in the literature \cite{chang2026}, was used as a
guide for setting the minimum area threshold $A_{min}$ in the pipeline.

\subsection{Defect Categories}\label{subsec:defect_categories}

The defect categories as shown in figure \ref{fig:dataset_preview} considered in this work include PbI$_2$ bright
particles, small pinholes, large pinholes, PbI$_2$ needle crystals,
and grain-boundary-associated defects. Pinholes and voids are defined
as local dark regions indicating incomplete film coverage. Bright
particles correspond to PbI$_2$ excess or undissolved crystallites
appearing as high-intensity spots. These regions are important because
they may contribute to charge leakage, non-radiative recombination,
reduced effective absorber coverage, and long-term instability in
perovskite photovoltaic devices \cite{chang2026,hu2023}.

\begin{table*}[!htbp]
\caption{Summary of the SEM image dataset used for defect quantification
and YOLO model training}
\label{tab:dataset_summary}
\small
\begin{tabular}{@{}lcccc@{}}
\toprule
Class / group & Base images & Patches (224$\times$224) & Defect type & Annotation \\
\midrule
PbI$_2$ excess (Class~0)          & 80  & \multirow{5}{*}{4{,}701 total} & Bright particles, needles & YOLO labels \\
3D pinholes (Class~1)              & 77  &                                 & Small/large dark pits     & YOLO labels \\
3D-2D pinholes (Class~2)           & 86  &                                 & Mixed pinholes            & YOLO labels \\
3D background (Class~3)            & 103 &                                 & Background texture        & YOLO labels \\
3D-2D background (Class~4)         & 94  &                                 & Background texture        & YOLO labels \\
\midrule
Total                              & 440 & 4{,}701                         & 5 classes                 & YOLO + manual \\
\botrule
\end{tabular}
\end{table*}

%---------------------------------------------------------------------
% Dataset preview figure
%---------------------------------------------------------------------
%\begin{figure*}[t]
%\centering
%\IfFileExists{dataset preview.png}{
%  \includegraphics[width=0.96\textwidth]{dataset preview.png}
%}
%\caption{Representative 224$\times$224~px SEM patches from the five morphology classes used %for YOLO training and MicroDefectCV evaluation.}
%\label{fig:dataset_preview}
%\end{figure*}
\begin{figure*}[!htbp]
\centering
\includegraphics[width=0.96\textwidth]{dataset preview.png}
\caption{Representative 224$\times$224~px SEM patches from the five morphology classes used for YOLO training and MicroDefectCV evaluation.}
\label{fig:dataset_preview}
\end{figure*}

\section{Proposed Methodology for Materials-Aware Defect Quantification}
\label{sec:methodology}

The proposed pipeline converts SEM morphology into quantitative defect
descriptors. It consists of image preprocessing, local contrast
enhancement, dynamic thresholding, morphological refinement, contour
extraction, and defect quantification. Figure~\ref{fig:pipeline} shows
the complete workflow.

\begin{figure*}[!htbp]
\centering
\IfFileExists{Fig1_pipeline_overview.png}
{\includegraphics[width=0.95\textwidth]{Fig1_pipeline_overview.png}}

\caption{MicroDefectCV detection pipeline: from raw SEM input through
  CLAHE enhancement, dynamic thresholding, and contour filtering
  to the final annotated defect mask.}
\label{fig:pipeline}
\end{figure*}

\subsection{Preprocessing and SEM Intensity Normalisation}
\label{subsec:preprocessing}

Each SEM image is converted to grayscale and normalised to reduce
image-to-image intensity variation. This step is necessary because
SEM images can differ in brightness, contrast, magnification, and
local charging effects. Noise reduction is applied before segmentation
to suppress isolated artefacts while preserving defect boundaries.

\subsection{CLAHE-Based Local Contrast Enhancement}
\label{subsec:clahe}

Contrast-limited adaptive histogram equalization (CLAHE) is applied
to enhance local intensity differences between defect and non-defect
regions. In SEM images of perovskite films, pinholes and voids
typically appear as dark regions with weak boundaries \cite{chang2026}.
CLAHE improves local contrast while limiting excessive noise
amplification, making subsequent thresholding more stable.

% CLAHE parameters are mode-dependent. Default values used across modes:
% clipLimit: 3.0; tileGridSize: (8, 8). See Table~\ref{tab:parameters}.

\subsection{Dynamic Thresholding for Candidate Defect Extraction}
\label{subsec:dynamic_thresholding}

Instead of using a fixed threshold for all SEM images, the proposed
method computes a dynamic threshold from image intensity statistics.
This improves adaptability when SEM images are captured under
different illumination or contrast conditions.

\begin{equation}
T = \mu - k\sigma
\label{eq:threshold}
\end{equation}

where $T$ is the dynamic threshold, $\mu$ is the mean pixel
intensity, $\sigma$ is the standard deviation of pixel intensities,
and $k$ is a sensitivity factor. Pixels darker than $T$ are selected
as candidate defect pixels:

\begin{equation}
M(x,y)=
\begin{cases}
1, & I(x,y)<T \\
0, & I(x,y)\geq T
\end{cases}
\label{eq:mask}
\end{equation}

where $I(x,y)$ is the intensity at pixel coordinate $(x,y)$, and
$M(x,y)$ is the binary candidate defect mask.

\subsection{Morphological Refinement}
\label{subsec:morphological_refinement}

The initial binary mask may contain small isolated noise regions or
fragmented defect boundaries. Morphological opening removes small
components, while morphological closing connects fragmented defect
regions. This step improves the structural consistency of the defect
mask before contour extraction.

\subsection{Contour-Based Filtering and Defect Quantification}
\label{subsec:contour_filtering}

Contours are extracted from the refined binary mask. Each candidate
region is filtered using material-relevant features such as area,
equivalent diameter, aspect ratio, circularity, and local intensity.
Very small components are removed to avoid treating SEM noise as
material defects. The minimum area threshold is informed by the
expected void diameter range of $\sim$100--200~nm reported for
similar perovskite films \cite{chang2026}.

The defect-area percentage is calculated as:

\begin{equation}
D_{area} = \frac{N_{defect}}{N_{total}} \times 100
\label{eq:defect_area}
\end{equation}

where $D_{area}$ is the defect-area percentage, $N_{defect}$ is the
number of pixels classified as defects, and $N_{total}$ is the total
number of pixels in the SEM image. Higher defect-area percentage may
indicate poorer film coverage or increased probability of leakage
pathways \cite{chang2026,hu2023}.

%=====================================================================
% 6. ALGORITHM AND IMPLEMENTATION
%=====================================================================

\subsection{Algorithm and Implementation}\label{sec:algorithm}

Algorithm~\ref{algo:pipeline} summarises the proposed SEM microdefect
quantification process.

\begin{algorithm}
\caption{SEM microdefect quantification for perovskite photovoltaic films}
\label{algo:pipeline}
\begin{algorithmic}[1]
\Function{QuantifyMicrodefects}{$I, k, A_{min}$}
    \Statex \qquad \textbullet\ \textbf{requirement 1}: SEM image $I$
    \Statex \qquad \textbullet\ \textbf{requirement 2}: sensitivity factor $k$ and minimum area $A_{min}$
    \Statex \qquad \textbullet\ \textbf{returns 1}: a binary defect mask $M$
    \Statex \qquad \textbullet\ \textbf{returns 2}: defect regions $R$ and quantitative descriptors $Q$
    \State $G \gets \Call{Grayscale}{I}$
    \State $G_n \gets \Call{NormalizeAndDenoise}{G}$
    \State $C \gets \Call{CLAHE}{G_n}$ \Comment{Local contrast enhancement}
    \State $\mu, \sigma \gets \Call{ComputeStatistics}{C}$
    \State $T \gets \mu - k\sigma$ \Comment{Dynamic threshold calculation}
    \State $M_c \gets \{(x,y) \mid C(x,y) < T\}$
    \State $M_c \gets \Call{MorphologicalFilter}{M_c}$ \Comment{Opening and closing operations}
    \State $\mathcal{C} \gets \Call{ExtractContours}{M_c}$
    \State $R \gets \emptyset$
    \For{every contour $c \in \mathcal{C}$}
        \If{$\Call{Area}{c} \ge A_{min}$}
            \State $R \gets R \cup \{c\}$
        \EndIf
    \EndFor
    \State $M \gets \Call{CreateMask}{R}$
    \State $Q \gets \Call{ComputeDescriptors}{R}$
    \State \textbf{return} $M, R, Q$
\EndFunction
\end{algorithmic}
\end{algorithm}

\subsubsection{The MicroDefectCV Package}\label{subsec:microdefectcv_pkg}

The pipeline described in this section is released as
\texttt{microdefectcv}~\cite{microdefectcv2025}, a pip-installable
Python package available at
\url{https://github.com/Sahilsonii/microdefectcv} and
\url{https://pypi.org/project/microdefectcv/}. The package provides
six detection modes tailored to different perovskite morphologies:
\texttt{pbi2} (bright crystallite particles),
\texttt{pinhole} (dark voids),
\texttt{2d} (flat 2D perovskite morphology),
\texttt{3d} (3D grains with boundary suppression),
\texttt{3d\_2d} (mixed morphology), and
\texttt{auto} (automatic mode that classifies the input image from
intensity statistics alone). Results are returned as binary masks,
bounding-box contours, defect count, and defect-area ratio, with
an optional YOLO-format label file export
(\texttt{save\_yolo\_annotations}). The package has a command-line
entry point (\texttt{microdefectcv}) and runs on a standard CPU
with no deep-learning dependency, achieving under one second per
image on an Intel Core~i7 system.

\subsubsection{Software and Reproducibility}\label{subsec:software}

The proposed method was implemented in Python~3.10 using OpenCV~4.8
and NumPy~1.24. The pipeline runs on a standard CPU and processes
a single SEM image in under one second on an Intel Core i7 system.
The YOLO training experiments were conducted on an NVIDIA RTX~3050~Ti
(4~GB VRAM) GPU using Ultralytics~8.0. Training was performed at
224$\times$224~pixel resolution for 50~epochs per model. The complete
code is available at the repository listed in the Code Availability
section.

\begin{table*}[!htbp]
\caption{Key parameter settings of the MicroDefectCV pipeline}
\label{tab:parameters}
\renewcommand{\arraystretch}{0.85}
\begin{tabular}{@{}llp{6cm}@{}}
\toprule
Parameter & Value & Purpose \\
\midrule
CLAHE clip limit          & 3.0        & Limits contrast amplification \\
CLAHE tile grid size      & $8\times8$ & Defines local histogram region \\
Gaussian blur $\sigma$    & 1.0        & Noise suppression before CLAHE \\
Dark percentile           & 12.0       & Candidate dark-region threshold \\
Strict dark percentile    & 6.0        & Fallback for high-fill images \\
Max fill ratio            & 0.08       & Guards against over-segmentation \\
Min area (dark), $A_{min}$& 15~px      & Removes sub-pixel noise \\
Max area (dark)           & 20{,}000~px& Excludes image-wide artefacts \\
Min circularity (dark)    & 0.30       & Shape filter for pinholes \\
Min solidity              & 0.40       & Compact region enforcement \\
NMS IoU threshold         & 0.40       & Removes overlapping boxes \\
Confidence threshold      & 0.30       & Minimum detection confidence \\
\botrule
\end{tabular}
\end{table*}

AS listing in table \ref{tab:parameters} the simplified thresholding
logic used for dark-region candidate extraction. The full
implementation in the \texttt{microdefectcv} package additionally
includes mode-specific grain-boundary suppression, top-hat
transforms for bright-particle detection, needle-crystal filtering,
and IoU-based non-maximum suppression.

%=====================================================================
% 7. EXPERIMENTAL VALIDATION
%=====================================================================

\section{Experimental Validation}\label{sec:validation}

The annotation and validation workflow proceeded through three
stages before arriving at the MicroDefectCV pipeline. First, two
automated zero-shot approaches (AutoDistill and SAM) were tested
and found insufficient for SEM morphology. Second, manual
bounding-box annotation was carried out to produce reliable
ground truth. Third, the MicroDefectCV classical pipeline was
developed and benchmarked against both the manual reference and
deep-learning YOLO models.

%---------------------------------------------------------------------
\subsection{Attempted Zero-Shot Auto-Annotation Using AutoDistill}
\label{subsec:autodistill}
%---------------------------------------------------------------------

Prior to committing to manual annotation, a zero-shot automatic
annotation approach was evaluated to accelerate ground-truth
generation. The framework selected was
AutoDistill~\cite{autodistill2023}, which composes Grounded
DINO~\cite{liu2023groundingdino} for text-prompted bounding-box
detection with SAM~\cite{kirillov2023sam} for pixel-level
segmentation. Ten natural-language prompts describing pinholes
(e.g., ``dark circular void on surface'', ``deep dark cavity'')
were submitted to the pipeline. Despite extensive prompt engineering,
AutoDistill exhibited \textbf{systematic under-detection in
high-density pinhole regions}: Grounding DINO's NMS suppressed
overlapping adjacent detections, the grayscale SEM domain conflicted
with the model's colour-image pre-training, and confidence scores
varied unpredictably with magnification. Figure~\ref{fig:autodistill_result}
shows the under-detection failure on a representative frame.

\begin{figure}[!htbp]
\centering
\IfFileExists{Fig_autodistill.png}
{\includegraphics[width=0.95\columnwidth]{Fig_autodistill.png}}
{
\fbox{\begin{minipage}[c][3.5cm][c]{0.90\columnwidth}\centering
\textbf{[Placeholder: AutoDistill output]}\\[4pt]
Upload comparison as \texttt{Fig\_autodistill.png}.
\end{minipage}}
}
\caption{AutoDistill output: only a fraction of pinholes are detected
  in dense regions due to NMS suppression and SEM domain gap.}
\label{fig:autodistill_result}
\end{figure}

%---------------------------------------------------------------------
\subsubsection{Alternative Attempt: Segment Anything Model (SAM)}
\label{subsubsec:sam_attempt}
%---------------------------------------------------------------------

To circumvent AutoDistill's recall limitations, the standalone
SAM \texttt{vit\_b} model was tested in automatic mask generation
mode with a low IoU threshold (\texttt{pred\_iou\_thresh=0.35}) and
downstream area filters targeting the pinhole size range. SAM
exhibited severe \textbf{over-segmentation}: it generated hundreds
of arbitrary grain-texture patches rather than isolating the dark
pinhole cavities. Because SAM is class-agnostic and relies on sharp
edges, and perovskite pinholes present as soft low-contrast gradients
in SEM, the masks did not correspond to material defect boundaries.
Figure~\ref{fig:sam_result} shows the over-segmentation failure.

\begin{figure}[!htbp]
\centering
\IfFileExists{Fig_sam_result.jpg}
{\includegraphics[width=0.95\columnwidth]{Fig_sam_result.jpg}}
{\fbox{\begin{minipage}[c][3.5cm][c]{0.90\columnwidth}\centering
\textbf{[Placeholder: SAM over-segmentation result]}\end{minipage}}}
\caption{SAM over-segments the SEM image into grain-texture patches
  rather than isolating pinholes.}
\label{fig:sam_result}
\end{figure}

%---------------------------------------------------------------------
\subsection{Manual Annotation Reference}\label{subsec:manual_reference}
%---------------------------------------------------------------------

Because both zero-shot approaches failed to produce reliable
ground truth, manual bounding-box annotation was adopted. Annotation
was performed using a custom OpenCV-based tool that displays each
SEM image on a canvas, allows annotators to draw bounding boxes by
class, and saves labels in YOLO format (\texttt{.txt} files with
normalised $x_c$, $y_c$, $w$, $h$ coordinates). PbI$_2$ excess
regions (Class~0) yielded more than 200~boxes per frame, while 3D
pinholes (Class~1) produced 5--15 sparse but large annotations per
frame. Figure~\ref{fig:anno_gt} shows representative annotations
for both classes.

\begin{figure*}[!htbp]
\centering
\begin{subfigure}[t]{0.48\textwidth}
  \centering
  \IfFileExists{Fig_anno_pbi2.png}{
    \includegraphics[width=\textwidth]{Fig_anno_pbi2.png}
  }{
    \fbox{\begin{minipage}[c][5cm][c]{\textwidth}\centering
    \texttt{Fig\_anno\_pbi2.png} not found.\end{minipage}}
  }
  \caption{PbI$_2$ excess (Class~0): dense orange bounding boxes.}
  \label{fig:anno_pbi2}
\end{subfigure}\hfill
\begin{subfigure}[t]{0.48\textwidth}
  \centering
  \IfFileExists{Fig_anno_pinhole.png}{
    \includegraphics[width=\textwidth]{Fig_anno_pinhole.png}
  }{
    \fbox{\begin{minipage}[c][5cm][c]{\textwidth}\centering
    \texttt{Fig\_anno\_pinhole.png} not found.\end{minipage}}
  }
  \caption{3D pinholes (Class~1): sparse cyan bounding boxes.}
  \label{fig:anno_pinhole}
\end{subfigure}
\caption{Manual bounding-box annotations for PbI$_2$ excess (a)
  and 3D pinhole (b) defect classes.}
\label{fig:anno_gt}
\end{figure*}

%---------------------------------------------------------------------
\subsubsection{Decision to Adopt Manual Annotation}
\label{subsubsec:manual_decision}
%---------------------------------------------------------------------

AutoDistill-generated labels carried a substantial false-negative
bias that any downstream YOLO model would inherit. SAM-generated
masks did not correspond to material defect boundaries. Manual
annotation was therefore the only approach that produced
ground-truth data reliable enough for YOLO model training and for
validating the MicroDefectCV classical pipeline.

%---------------------------------------------------------------------
\subsection{MicroDefectCV Adaptive Filter Pipeline}
\label{subsec:baselines}
%---------------------------------------------------------------------

Following the annotation phase, the MicroDefectCV adaptive filter
pipeline was applied to the same SEM images. Unlike the annotation
approaches above, MicroDefectCV requires no ground-truth labels.
It selects a detection mode automatically from image statistics
(\texttt{auto} mode) and runs CLAHE, dynamic thresholding,
grain-boundary suppression, and contour shape filtering to produce
binary defect masks and YOLO-format annotations in under one second
per image on CPU. The pipeline output was compared against the
manual annotation ground truth and the following classical baselines:

\begin{itemize}
\item Otsu thresholding
\item CLAHE followed by Otsu thresholding
\item Canny edge detection
\item Watershed segmentation
\end{itemize}

A visual qualitative comparison of these methods against the MicroDefectCV pipeline and manual annotations is shown in Figure~\ref{fig:qualitative} (Section~\ref{sec:results}).

%=====================================================================
% 8. RESULTS
%=====================================================================

\section{Results and Discussion}\label{sec:results}
The proposed pipeline generated interpretable defect masks and
localised pinhole-like regions in SEM images. Qualitative comparison
showed that the proposed method produced more complete defect regions
than Canny edge detection and reduced over-segmentation compared
with watershed segmentation. Figure ~\ref{fig:qualitative} shows
representative SEM defect detection results.
\begin{figure*}[!htbp]
\centering
\IfFileExists{Fig2_sem_pipeline.png.png}
{\includegraphics[width=0.95\textwidth]{Fig2_sem_pipeline.png.png}}

\caption{Qualitative comparison: MicroDefectCV versus manual annotation
  and classical baselines on representative perovskite SEM patches.}
\label{fig:qualitative}
\end{figure*}

\subsection{Quantitative Segmentation Performance}
\label{subsec:quantitative}

Table~\ref{tab:segmentation_metrics} compares pixel-level precision,
recall, F1-score, and IoU for MicroDefectCV against classical
baselines, evaluated against the manual annotation ground truth.
MicroDefectCV's mode-aware design reduces false positives from
grain-boundary texture, which commonly afflicts global Otsu
thresholding. Canny edge detection returns boundary maps rather
than filled defect regions, making area-based quantification
inapplicable. Watershed segmentation tends to over-segment in
noisy perovskite SEM texture. The dynamic percentile threshold in
MicroDefectCV adapts to per-image contrast variation, providing
more consistent recall across different morphology classes.

\begin{table*}[!htbp]
\centering
\caption{Segmentation performance of MicroDefectCV versus classical baselines}
\label{tab:segmentation_metrics}
\begin{tabular}{@{}lcccc@{}}
\toprule
Method & Precision & Recall & F1-score & IoU \\
\midrule
Otsu thresholding & 0.089 & 0.379 & 0.144 & 0.078 \\
CLAHE + Otsu & 0.057 & 0.58 & 0.104 & 0.055 \\
Canny edge detection & 0.053 & 0.649 & 0.098 & 0.052 \\
Adaptive thresholding & 0.063 & 0.42 & 0.11 & 0.058 \\
\textbf{MicroDefectCV} & \textbf{0.442} & \textbf{0.127} & \textbf{0.197} & \textbf{0.109} \\
\botrule
\end{tabular}
\end{table*}

\subsection{Defect Quantification Results}
\label{subsec:defect_quantification_results}

\begin{table*}[!htbp]
\caption{Material-relevant defect descriptors extracted from SEM images by MicroDefectCV}
\label{tab:defect_descriptors}
\begin{tabular}{@{}lccccc@{}}
\toprule
Sample & Defect count & Total area (px) & Mean diameter (nm) & Defect-area (\%) \\
\midrule
Sample 1 & 4 & 304 & 4.9 & 0.46 \\
Sample 2 & 28 & 1692 & 4.4 & 2.58 \\
Sample 3 & 1 & 170 & 7.4 & 0.26 \\
Sample 4 & 28 & 2288 & 5.1 & 3.49 \\
\botrule
\end{tabular}
\end{table*}

Table~\ref{tab:defect_descriptors} shows the type of morphology
descriptors produced by the proposed pipeline. A mean defect diameter
in the range of $\sim$100--200~nm would be consistent with the
nanovoid dimensions reported by Chang et al.\ \cite{chang2026} for
similar perovskite film systems, providing a materials-grounded
interpretation of the pipeline output.

\subsection{Deep-Learning Defect Detection Results}
\label{subsec:yolo_results}

To benchmark the classical pipeline against a data-driven approach,
six YOLO models were trained on the 4{,}701-patch dataset using
the five-class annotation scheme described in
Section~\ref{subsec:sem_dataset}. Models from both the YOLOv8 \cite{ref23} and
YOLO11 families were evaluated, spanning small (s), medium (m), and
large (l) backbone configurations. All models were trained for
50~epochs at 224$\times$224~pixel resolution on an NVIDIA RTX~3050~Ti
(4~GB VRAM) GPU using Ultralytics 8.0 \cite{ultralytics2023}.

Table~\ref{tab:yolo_results} reports the precision, recall,
F1-score, mAP50, mAP50-95, and training time for each model. The
architecture of the top-performing YOLO11m model is illustrated in
Figure~\ref{fig:yolo_arch}, highlighting the C3k2-based backbone
optimised for multi-scale feature extraction.

\begin{figure*}[!htbp]
\centering
\IfFileExists{Fig3_yolo11_arch.png}{
  \includegraphics[width=0.96\textwidth]{Fig3_yolo11_arch.png}
}{
  \fbox{\begin{minipage}[c][6cm][c]{0.92\textwidth}\centering
  Place \texttt{Fig3\_yolo11\_arch.png} in the directory.\end{minipage}}
}
\caption{YOLO11 architecture used for SEM defect detection benchmarking, showing the C3k2-backbone, SPPF, and detection head configuration.}
\label{fig:yolo_arch}
\end{figure*}

\begin{table*}[!htbp]
\centering
\caption{YOLO model comparison on the 4{,}701-patch perovskite SEM dataset
(50~epochs, 224$\times$224~px, NVIDIA RTX~3050~Ti 4~GB VRAM)}
\label{tab:yolo_results}
\begin{tabular}{@{}lcccccc@{}}
\toprule
Model & Precision & Recall & F1-score & mAP50 & mAP50-95 & Train time \\
\midrule
YOLOv8s  & 0.455 & 0.425 & 0.440 & 0.378 & 0.163 & 27.7 min \\
YOLOv8m  & 0.398 & 0.430 & 0.413 & 0.353 & 0.154 & 39.1 min \\
YOLOv8l  & 0.304 & 0.308 & 0.306 & 0.237 & 0.089 & 70.3 min \\
YOLO11s  & 0.433 & 0.458 & 0.445 & 0.396 & 0.173 & 37.3 min \\
\textbf{YOLO11m} & \textbf{0.449} & \textbf{0.482} & \textbf{0.465} & \textbf{0.422} & \textbf{0.191} & 68.5 min \\
YOLO11l  & 0.405 & 0.432 & 0.418 & 0.353 & 0.160 & \textbf{176.8 min} \\
\botrule
\end{tabular}
\end{table*}

YOLO11m achieved the best performance (mAP50~=~0.422,
mAP50-95~=~0.191, precision~=~0.449, recall~=~0.482). Despite this,
it required 68.5~min of GPU training on 4{,}701 labelled patches,
contrasted with MicroDefectCV which detects defects in under one
second per image with zero labelling effort. Training wall-clock
times for all six models are summarised in
Figure \ref{fig:training_time}; YOLO11l required 176.8~min while
providing no additional accuracy gain over YOLO11m, confirming that
dataset size rather than model capacity is the main bottleneck.

\begin{figure}[!htbp]
\centering
\IfFileExists{training_time.png}{
  \includegraphics[width=0.70\columnwidth]{training_time.png}
}{
  \fbox{\begin{minipage}[c][3cm][c]{0.65\columnwidth}\centering
  \texttt{training\_time.png} not found.\end{minipage}}
}
\caption{Training wall-clock time for all six YOLO variants.}
\label{fig:training_time}
\end{figure}

\begin{figure*}[!htbp]
\centering
\begin{subfigure}[t]{0.48\textwidth}
  \centering
  \IfFileExists{model_map5095_comparison.png}{
    \includegraphics[width=\textwidth]{model_map5095_comparison.png}
  }{
    \fbox{\begin{minipage}[c][5cm][c]{\textwidth}\centering
    \texttt{model\_map5095\_comparison.png} not found.\end{minipage}}
  }
  \caption{mAP50-95 comparison: YOLO11m achieves the highest score (0.191) among all six variants.}
  \label{fig:map_comparison}
\end{subfigure}\hfill
\begin{subfigure}[t]{0.48\textwidth}
  \centering
  \IfFileExists{precision_recall_scatter.png}{
    \includegraphics[width=\textwidth]{precision_recall_scatter.png}
  }{
    \fbox{\begin{minipage}[c][5cm][c]{\textwidth}\centering
    \texttt{precision\_recall\_scatter.png} not found.\end{minipage}}
  }
  \caption{Precision-recall scatter: YOLO11m (gold) achieves precision 0.449 / recall 0.482, best among all variants.}
  \label{fig:pr_scatter}
\end{subfigure}
\caption{mAP50-95 bar chart (a) and precision-recall scatter (b)
  for all six YOLO variants; YOLO11m dominates both panels.}
\label{fig:model_selection}
\end{figure*}

As shown in figure \ref{fig:model_selection} The moderate mAP50-95 values across all models (0.089--0.191) reflect
the inherent difficulty of detecting nanoscale microdefects in SEM
images. Perovskite pinholes and voids appear as small, weakly
contrasted dark regions at 224$\times$224~pixel resolution, making
high-precision bounding-box localisation at strict IoU thresholds
challenging. As depict in figure \ref{fig:training_curves} and figure \ref{fig:yolo_eval} these results are consistent with the known difficulty
of SEM defect detection at nanoscale resolution and highlight the
current dataset size as the main constraint on model performance.
The trained YOLO11m weights are available at the repository listed
in the Code Availability section.

% ── YOLO11m Training Curves ──
\begin{figure*}[!htbp]
\centering
\IfFileExists{Fig6_training_curves.png}{
  \includegraphics[width=0.96\textwidth]{Fig6_training_curves.png}
}{
  \fbox{\begin{minipage}[c][4cm][c]{0.92\textwidth}\centering
  Place \texttt{Fig6\_training\_curves.png} in the directory.\end{minipage}}
}
\caption{YOLO11m training and validation loss/mAP curves over
  50 epochs, confirming stable convergence within the 4~GB VRAM budget.}
\label{fig:training_curves}
\end{figure*}

% ── Confusion Matrix + PR Curve ──
%---------------------

\begin{figure*}[!htbp]
\centering

\begin{subfigure}[t]{0.54\textwidth}
  \centering
  \IfFileExists{confusion_matrix_normalized.png}{
    \includegraphics[width=\textwidth]{confusion_matrix_normalized.png}
  }
  \caption{}
  \label{fig:confusion_matrix}
\end{subfigure}
\hfill
\begin{subfigure}[t]{0.42\textwidth}
  \centering
  \IfFileExists{BoxPR_curve.png}{
    \includegraphics[width=\textwidth]{BoxPR_curve.png}
  }
  \caption{}
  \label{fig:pr_curve}
\end{subfigure}

\caption{YOLO11m evaluation: (a) normalised confusion matrix; 
(b) per-class precision-recall curve on the held-out test set.}
\label{fig:yolo_eval}
\end{figure*}

% ── Label Distribution + Validation Predictions ──
\begin{figure*}[!htbp]
\centering
  \includegraphics[width=0.92\textwidth]{Pbl2.png}\\[4pt]
  \includegraphics[width=0.92\textwidth]{3D-pinholes.png}\\[4pt]
  \includegraphics[width=0.92\textwidth]{3D-2D pinholes.png}
  \caption{Representative YOLO11m predictions for PbI$_2$ excess, 3D pinholes, and 3D-2D pinholes (from top to bottom).}
  \label{fig:val_pred}
\label{fig:yolo_data}
\end{figure*}

%=====================================================================
% 9. DISCUSSION
%=====================================================================

\subsection{Discussion}\label{sec:discussion}

The results shown in figure \ref{fig:yolo_data} suggest that automated SEM morphology quantification can support materials characterisation of perovskite photovoltaic films.
The proposed method is not limited to visual enhancement; it produces
measurable descriptors that can be linked to film coverage, defect
density, and potential device-quality issues.

\subsubsection{Physical Interpretation of Detected Defects}

The dark regions detected by the proposed pipeline correspond to
pinhole-like and void-like morphological features. Based on recent
experimental literature, these features carry the following physical
significance. Voids of the size range targeted by the pipeline
($\sim$100--200~nm equivalent diameter) are consistent with those
reported by Chang et al.\ \cite{chang2026} in 1.67~eV wide-bandgap
perovskite films. These voids were experimentally linked to increased
trap-state density ($2.62 \times 10^{16}$~cm$^{-3}$) and reduced fill
factor. A pipeline that can automatically detect and measure such
features provides a quantitative route to assess whether a given
fabrication condition is likely to produce trap-rich or trap-lean
film morphology.

The defect-area percentage (Equation~\ref{eq:defect_area}) can be
interpreted as a proxy for the fraction of the absorber area that is
morphologically compromised. Higher defect-area percentage may
correspond to greater non-radiative recombination losses, reduced
effective absorber coverage, and increased probability of shunting
pathways. In flexible and fiber-related devices, spatial gradients in
defect-area percentage along the device geometry may indicate local
coating non-uniformities caused by substrate curvature or tension.

\subsubsection{Comparison with Baseline Methods}

CLAHE improved the visibility of weakly contrasted pinhole-like
regions, while dynamic thresholding reduced dependence on a single
manually chosen threshold. Morphological refinement improved the
continuity of defect masks. Contour-based filtering reduced false
positives by excluding regions that were too small or geometrically
inconsistent with material defects.

Compared with Otsu thresholding, the proposed approach provides more
direct control over defect sensitivity through the parameter $k$.
Compared with Canny edge detection, it provides filled defect regions
rather than boundary-only maps, enabling area-based quantification.
Compared with watershed segmentation, it reduces over-segmentation
in noisy SEM textures.

\subsubsection{Connection to Manufacturing-Scale Quality Control}

The transition from laboratory-scale to manufacturing-scale perovskite
photovoltaics requires rapid, reproducible feedback between film
deposition and morphology quality. As demonstrated in recent work on
large-area triple-junction devices \cite{zheng2025}, SEM imaging is
used to confirm that performance benchmarks are maintained when
scaling from 1~cm$^2$ to 16~cm$^2$ device areas. Automated morphology
quantification pipelines of the type proposed here are a necessary
step toward making this morphology feedback loop scalable and
objective.

Future work should validate the correlation between pipeline-derived
defect descriptors and device performance parameters such as fill
factor, open-circuit voltage, and stability---following the
experimental framework used by Chang et al.\ \cite{chang2026}, who
showed that a 14.8\% reduction in trap-state density corresponds to
measurable improvements in V$_{oc}$, FF, and PCE.

%\subsection{Limitations}
%The method may still classify naturally dark texture regions as
%defects when they resemble pinholes. Very small defects below the
%minimum area threshold may be missed. The method also depends on
%image quality, magnification, and SEM acquisition settings. The
%proposed pipeline should therefore be viewed as a reproducible
%screening and quantification tool rather than a complete replacement
%for expert materials interpretation.

%=====================================================================
% 10. ENERGY-DEVICE AND MANUFACTURING RELEVANCE
%=====================================================================

\subsection{Energy-Device and Manufacturing Relevance}
\label{sec:energy_relevance}

Scalable perovskite photovoltaic manufacturing requires rapid feedback
between material processing and device-quality assessment. The
importance of this feedback loop is evident from recent high-performance
demonstrations: Chang et al.\ \cite{chang2026} showed that optimising
hole-transport layer morphology to eliminate nanovoids improved the
certified efficiency of 1.67~eV wide-bandgap perovskite cells from
21.92\% to 23.42\%, with the fill factor increasing to 86.8\%. This
level of improvement, achieved through a morphological change
detectable in SEM images, demonstrates the direct manufacturing value
of SEM-based defect quantification.

At the system level, morphological quality of individual subcell
films has been shown to be a key factor in achieving high efficiency
in tandem \cite{chang2026} and triple-junction \cite{zheng2025}
architectures. In both types of devices, SEM imaging was used to
confirm the absence of voids and delamination before performance
testing, indicating that morphological verification is a standard
quality gate.

For fiber-related and flexible perovskite photovoltaic devices, the
proposed method can support quality assessment along curved surfaces
or mechanically deformable substrates. Defect descriptors such as
defect density and defect-area percentage may be used to identify
coating non-uniformity, local film failure, or morphology degradation
after bending tests. In future work, these descriptors should be
correlated with photovoltaic parameters such as open-circuit voltage,
short-circuit current density, fill factor, power conversion
efficiency, and stability retention.

%=====================================================================
% 11. LIMITATIONS AND FUTURE WORK
%=====================================================================

\section{Limitations and Future Work}\label{sec:limitations}

This study has several limitations. First, the method depends on SEM
image quality and may require parameter adjustment for different
magnifications or imaging conditions. Second, the current pipeline
is designed mainly for dark pinhole-like and void-like defects; other
defect types such as bright particles, cracks, or delamination
regions may require additional features. Third, the method should be
validated on larger datasets containing different perovskite
compositions, coating methods, flexible substrates, and fiber-shaped
device geometries. Fourth, future work should correlate extracted
defect descriptors with photovoltaic performance and stability data,
following the experimental framework established by Chang et al.\
\cite{chang2026}.



%=====================================================================
% 12. CONCLUSION
%=====================================================================

\section*{Conclusion and Future Work}\label{sec:conclusion}

This work presented MicroDefectCV, an open-source, pip-installable
Python library for adaptive OpenCV-based defect detection and
quantification in perovskite photovoltaic SEM images~\cite{microdefectcv2025}.
The package implements six morphology-aware detection modes and an
automatic image-type classifier, requiring no deep learning, labelled
data, or GPU resources, and processing a single SEM image in under
one second on CPU. It exports results directly in YOLO bounding-box
format, eliminating the annotation bottleneck that otherwise requires
hundreds of hours of expert labelling. To contextualise the classical
approach, three annotation strategies were evaluated (zero-shot
AutoDistill, manual annotation, and standalone SAM), and six YOLO
models were benchmarked on the resulting 4{,}701-patch, five-class
dataset. The best deep-learning model (YOLO11m, mAP50~=~0.422)
outperformed MicroDefectCV on detection precision metrics, but
required 68.5~min of GPU training and thousands of labelled bounding
boxes --- an overhead that MicroDefectCV entirely avoids. The detected
defect size range ($\sim$100--200~nm equivalent diameter) is consistent
with nanovoids experiment where such features directly increase trap-state density and reduce fill factor. 
%ally documented by Chang et al.~\cite{chang2026},

Future extensions may include deep learning-assisted segmentation,
YOLO-based defect localisation, uncertainty-aware defect
classification, and integration with high-throughput fabrication
workflows. The proposed interpretable pipeline can also serve as a
pseudo-labelling or preprocessing tool for building larger annotated
SEM datasets. Future work will validate the pipeline on larger datasets, extend it to curved and fiber-shaped device geometries, and correlate SEM defect descriptors
with photovoltaic performance data.

%=====================================================================
% BACK MATTER
%=====================================================================

\backmatter

%\bmhead{Supplementary information}
%No supplementary information is available for this manuscript.

%\bmhead{Acknowledgements}
%The author acknowledges the institutional and laboratory support
%received during this work.

% TODO: Replace with exact acknowledgement names and departments.

\clearpage
\section*{Statements and Declarations}

\smallskip
\noindent\textbf{Funding}

No funding was received to assist with the preparation of this manuscript.

\smallskip
\noindent\textbf{Competing interests}

The author declares that there are no competing interests relevant to the content of this article.

\smallskip
\noindent\textbf{Ethics approval and consent to participate}

Not applicable. This study does not involve human participants or animals.

\smallskip
\noindent\textbf{Consent for publication}

Not applicable.

\smallskip
\noindent\textbf{Data availability}

The datasets generated and/or analysed during the current study are available from the corresponding author on reasonable request.

\smallskip
\noindent\textbf{Materials availability}

Not applicable.

\smallskip
\noindent\textbf{Code availability}

Experimental code, dataset preparation scripts, YOLO training pipelines, and the Streamlit annotation interface are available at \url{https://github.com/Sahilsonii/SEM-Based-Microdefect-Quantification-for-Fiber-Related-and-Flexible-Perovskite-Photovoltaic-Material}.

The \textbf{MicroDefectCV} classical pipeline is available as a pip-installable Python package (\texttt{pip install microdefectcv}) at \url{https://github.com/Sahilsonii/microdefectcv} and \url{https://pypi.org/project/microdefectcv/}.

\smallskip
\noindent\textbf{Author contribution}

Sahil Soni contributed to conceptualization, methodology, software implementation, experiments, analysis, visualization, and manuscript preparation. Dr. Archana Y. Chaudhari reviewed and approved the final manuscript.

\smallskip
\noindent\textbf{Use of AI-based tools}

AI-based tools were used only for language refinement, structural organisation, and LaTeX formatting assistance. The author verified the technical content, experimental results, references, and final manuscript.

%=====================================================================
% APPENDIX
%=====================================================================

\begin{appendices}
% Appendix A removed.
\end{appendices}

%=====================================================================
% REFERENCES
%=====================================================================

\begin{thebibliography}{99}

\bibitem{kojima2009}
Kojima A, Teshima K, Shirai Y, Miyasaka T.
Organometal halide perovskites as visible-light sensitizers for
photovoltaic cells.
J Am Chem Soc. 2009;131:6050--6051.
\url{https://doi.org/10.1021/ja809598r}

\bibitem{snaith2013}
Snaith HJ.
Perovskites: the emergence of a new era for low-cost, high-efficiency
solar cells.
J Phys Chem Lett. 2013;4:3623--3630.
\url{https://doi.org/10.1021/jz4020162}

\bibitem{green2014}
Green MA, Ho-Baillie A, Snaith HJ.
The emergence of perovskite solar cells.
Nat Photonics. 2014;8:506--514.
\url{https://doi.org/10.1038/nphoton.2014.134}

\bibitem{park2015}
Park NG.
Perovskite solar cells: an emerging photovoltaic technology.
Mater Today. 2015;18:65--72.
\url{https://doi.org/10.1016/j.mattod.2014.07.007}

\bibitem{green2025}
Green MA, Dunlop ED, Yoshita M, Kopidakis N, Bothe K, et al.
Solar cell efficiency tables (version~66).
Prog Photovolt Res Appl. 2025;33:795--810.
\url{https://doi.org/10.1002/pip.3919}

\bibitem{chang2026}
Chang L-C, Duong T, Ahmad V, Zhan H, Bui AD, et al.
High-efficiency perovskite/silicon tandem solar cells based on
wide-bandgap perovskite solar cells with unprecedented fill factor.
Nano-Micro Lett. 2026;18:122.
\url{https://doi.org/10.1007/s40820-025-01959-y}

\bibitem{zheng2025}
Zheng J, Wang G, Duan L, Duan W, Jiang Y, et al.
Tailoring nanoscale interfaces for perovskite--perovskite--silicon
triple-junction solar cells.
Nat Nanotechnol. 2025;20:1648--1655.
\url{https://doi.org/10.1038/s41565-025-02015-x}

\bibitem{balilonda2019}
Balilonda A, Li Q, Zhu M.
Perovskite solar fibers: current status, issues and challenges.
Adv Fiber Mater. 2019;1:101--125.
\url{https://doi.org/10.1007/s42765-019-00011-0}

\bibitem{hu2016}
Hu H, Yan K, Peng M, Yu X, Chen S, Chen B, et al.
Fiber-shaped perovskite solar cells with 5.3\% efficiency.
J Mater Chem A. 2016;4:3901--3906.
\url{https://doi.org/10.1039/C5TA10685A}

\bibitem{qiu2016}
Qiu L, He S, Yang J, Deng J, Peng H.
Fiber-shaped perovskite solar cells with high power conversion efficiency.
Small. 2016;12:2419--2424.
\url{https://doi.org/10.1002/smll.201600104}

\bibitem{diGiacomo2016}
Di Giacomo F, Fakharuddin A, Jose R, Brown TM.
Progress, challenges and perspectives in flexible perovskite solar cells.
Energy Environ Sci. 2016;9:3007--3035.
\url{https://doi.org/10.1039/C6EE01137C}

\bibitem{hu2023}
Hu H, Ritzer DB, Diercks A, Li Y, Singh R, et al.
Void-free buried interface for scalable processing of p-i-n-based
FAPbI\textsubscript{3} perovskite solar modules.
Joule. 2023;7:1574--1592.
\url{https://doi.org/10.1016/j.joule.2023.05.017}

\bibitem{chen2021}
Chen S, Dai X, Xu S, Jiao H, Zhao L, Huang J.
Stabilizing perovskite--substrate interfaces for high-performance
perovskite modules.
Science. 2021;373:902--907.
\url{https://doi.org/10.1126/science.abi6323}

\bibitem{alashouri2020}
Al-Ashouri A, K\"{o}hnen E, Li B, Magomedov A, Hempel H, et al.
Monolithic perovskite/silicon tandem solar cell with $>$29\% efficiency
by enhanced hole extraction.
Science. 2020;370:1300--1309.
\url{https://doi.org/10.1126/science.abd4016}

\bibitem{otsu1979}
Otsu N.
A threshold selection method from gray-level histograms.
IEEE Trans Syst Man Cybern. 1979;9:62--66.
\url{https://doi.org/10.1109/TSMC.1979.4310076}

\bibitem{canny1986}
Canny J.
A computational approach to edge detection.
IEEE Trans Pattern Anal Mach Intell. 1986;PAMI-8:679--698.
\url{https://doi.org/10.1109/TPAMI.1986.4767851}

\bibitem{vincent1991}
Vincent L, Soille P.
Watersheds in digital spaces: an efficient algorithm based on
immersion simulations.
IEEE Trans Pattern Anal Mach Intell. 1991;13:583--598.
\url{https://doi.org/10.1109/34.87344}

\bibitem{zuiderveld1994}
Zuiderveld K.
Contrast limited adaptive histogram equalization.
In: Heckbert PS, editor. Graphics Gems IV.
San Diego: Academic Press; 1994. p. 474--485.
\url{https://doi.org/10.1016/B978-0-12-336156-1.50061-6}

\bibitem{ronneberger2015}
Ronneberger O, Fischer P, Brox T.
U-Net: convolutional networks for biomedical image segmentation.
In: Medical Image Computing and Computer-Assisted Intervention --
MICCAI 2015.
Cham: Springer; 2015. p. 234--241.
\url{https://doi.org/10.1007/978-3-319-24574-4_28}

\bibitem{redmon2016}
Redmon J, Divvala S, Girshick R, Farhadi A.
You only look once: unified, real-time object detection.
In: Proceedings of the IEEE Conference on Computer Vision and
Pattern Recognition. 2016. p. 779--788.
\url{https://doi.org/10.1109/CVPR.2016.91}

\bibitem{ultralytics2023}
Jocher G, Chaurasia A, Qiu J.
Ultralytics YOLO. Version 8.0.0.
2023. \url{https://github.com/ultralytics/ultralytics}

\bibitem{autodistill2023}
Roboflow Inc.
AutoDistill: Label datasets with foundation models.
Version 0.1.
2023. \url{https://github.com/autodistill/autodistill}

\bibitem{liu2023groundingdino}
Liu S, Zeng Z, Ren T, Li F, Zhang H, Yang J, et al.
Grounding DINO: Marrying DINO with grounded pre-training for
open-set object detection.
arXiv preprint. 2023.
\url{https://doi.org/10.48550/arXiv.2303.05499}

\bibitem{microdefectcv2025}
Soni S.
MicroDefectCV: Adaptive OpenCV-based defect segmentation for
SEM and microstructure images.
Version 1.0. 2025.
\url{https://github.com/Sahilsonii/microdefectcv}
PyPI: \url{https://pypi.org/project/microdefectcv/}

\bibitem{kirillov2023sam}
Kirillov A, Mintun E, Ravi N, Mao H, Rolland C, Gustafson L, et al.
Segment anything.
In: Proceedings of the IEEE/CVF International Conference on Computer
Vision (ICCV). 2023. p. 4015--4026.
\url{https://doi.org/10.1109/ICCV51070.2023.00371}

\bibitem{ref23}
% TODO: Replace with the correct reference for the FAPbI3 one-pot
% synthesis and 2D/3D surface treatment protocol (GreatCell Solar materials).
Archana Y. Chaudhari, ,Prajwal Birwadkar, Sagar Joshi, Yash Verma, Rutuja Sindgi.
Classification of periapical dental X-ray using the YOLOv8 deep learning model.
MethodsX. 2025;15:103721.
\url{https://doi.org/10.1016/j.mex.2025.103721}

\bibitem{ref24}
Ansari, Z.A., Soni, S., Fatima, S. et al. A multi-model deep learning framework for SEM-based defect detection in FAPbI3 Perovskite thin films.
Scientific Report. 2025;15:41909.
\url{https://doi.org/10.1038/s41598-025-25848-x}
\end{thebibliography}

\end{document}